\chapter{Implementation}
\label{ch:implementation}

\section{Layout and state}

Particle state lives on the device for the whole run as thirteen float arrays:
three positions, three staged positions for the fused path, three velocities,
three accelerations and the mass. At two hundred thousand particles that is
about ten megabytes, so memory is never a constraint at these sizes and the
layout decision is entirely about access pattern.

The host holds no particle data during a run. Configuration and diagnostic
accumulators are the only host allocations, and the snapshot staging buffer is
touched only when a snapshot is actually written.

\section{Separating the physics from the kernels}

The force law, the smoothing kernels, the equation of state and the spatial
hash all live in headers marked so they compile both as device code and as
ordinary host C++. This buys two things. The kernels and the host side tests
cannot drift apart, because they are the same code. And the formula level
tests, which are the majority of them, run in continuous integration on a
machine with no GPU.

\section{The three force kernels}

The tiled kernel has one subtlety worth recording. Threads whose index is past
the end of the array must not return early, because they still have to reach
both barriers on every iteration of the tile loop; a barrier that only some
threads of a block arrive at is undefined behaviour. They run the whole loop
with clamped reads and discard the result at the end. Padded tile slots carry
zero mass, so they contribute nothing and the inner loop needs no bounds test
at all. Self interaction is likewise harmless and is deliberately not
special-cased: the numerator vanishes while the softened denominator stays
finite, so the diagonal term contributes exactly zero and the inner loop stays
branch free.

\section{The fluid step}

One fluid step is six device passes: hash, sort, cell bounds, gather, density,
and a combined force and integration pass.

The sort is a CUB radix sort of bucket and index pairs \cite{cccl}. Only the
bits that can differ are sorted: the table size bounds the bucket index, so the
sort stops at the significant bit count rather than walking all thirty two
bits.

After sorting, particle data is gathered into cell sorted order, and the
permutation is kept. The particle arrays end a step in the order the next step
wants, so the sort has an effect that outlives the step that paid for it.
Because that permutation must be a bijection, one gate checks that the multiset
of particle masses is unchanged across a step and another checks that total mass
is conserved over two thousand steps; a gather that lost or duplicated an index
would show up in both.

The force and integration pass is fused for the same reason the gravitational
one is, and it is safe for the same reason: the sorted arrays it reads are a
separate copy from the primary arrays it writes.

\section{Diagnostics without perturbing the run}

The particle state is single precision, but the reduction that measures
conservation is not. A float32 tree reduction over ten thousand terms loses
three to four decimal digits, and the energy gate resolves one part in a
thousand over ten thousand steps; reducing in single precision would have
measured the reduction rather than the integrator. Accumulators are double, and
the block partials are summed in a fixed order on the host rather than with
atomics, so two runs of the same configuration agree bit for bit.

The fused variant carries the leapfrog half step velocity, which is not the
velocity belonging to the current positions. Recovering the synchronised
velocity by shifting the array forward and then back would not return the
original bits, and that would break the bit identity guarantee for any run that
happened to take a diagnostic. The correction is applied inside the reduction
kernel instead, and a test asserts that taking a diagnostic on every step
leaves the trajectory unchanged.

\section{Snapshots and resume}

The snapshot header is sixteen bytes: magic number, particle count, step index,
and a flags word. The flags word carries the one thing a reader cannot infer
from the numbers, namely whether the stored velocities are synchronised with
the positions or offset by half a step. Resuming a fused run under the wrong
convention would shift the whole trajectory by half a step, which looks exactly
like a subtly wrong force.

Accelerations are not stored. They are a deterministic function of the
positions, so recomputing them on load is cheaper than storing them and
guarantees the resumed run enters the loop in exactly the state the original
run held.

\section{The style gate}

The repository forbids en dashes and em dashes. Enforcing that on source files
is easy; enforcing it on the reports is not, because TeX silently turns two
hyphens into an en dash. The gate therefore also parses the LaTeX sources for
that ligature, excluding verbatim regions where a command line flag legitimately
contains two hyphens, and runs over the compiled PDFs through a text extractor.
It caught a real instance during this build: the page ranges in the
bibliography, written in the conventional BibTeX style, would have rendered as
en dashes in the reference list of the finished document.
